% concatenated from main.tex
\documentclass[12pt,reqno]{amsart}
\usepackage[T1]{fontenc}
\usepackage{lmodern}
\usepackage[margin=1.2in]{geometry}
\usepackage{mathrsfs}
\usepackage[hidelinks]{hyperref}
\hypersetup{pdfstartview={FitH}}

\raggedbottom

\newtheorem{theorem}{Theorem}
\newtheorem{lemma}{Lemma}[section]
\newtheorem{proposition}[lemma]{Proposition}
\theoremstyle{remark}
\newtheorem{remark}{Remark}
\theoremstyle{definition}
\newtheorem{problem}{Problem}

\numberwithin{equation}{section}

\title[A 21-coloring of the plane]{A 21-Coloring of the Plane Without Monochromatic Unit-Area Rectangles}
\author[Gary Hu]{Gary Hu}
\address{Williams College}
\email{gh7@williams.edu}
\author[Yumo Liu]{Yumo Liu}
\address{Bard College}
\email{yy2323@bard.edu}
\date{}

\subjclass[2020]{Primary 05D10; Secondary 52C10}
\keywords{Euclidean Ramsey theory, plane coloring, rectangles, triangular lattice, hexagonal tiling}

\begin{document}

\begin{abstract}
Erd\H{o}s and Graham asked whether every finite coloring of the plane must contain a monochromatic rectangle of any prescribed area. Kova\v{c} gave a negative answer by constructing a $25$-coloring with no monochromatic rectangle of area $1$. We reduce the number of colors to $21$ by replacing the square cells in his construction with regular hexagons.
\end{abstract}

\maketitle

\section{Introduction}

Euclidean Ramsey theory asks which finite point configurations must occur monochromatically in every finite coloring of Euclidean space. In 1980, Graham proved that every finite coloring of $\mathbb R^n$ contains a monochromatic $n$-simplex with any prescribed positive volume \cite{graham}. Erd\H{o}s and Graham asked whether an analogous statement holds for rectangles and parallelograms \cite{erdos-graham}. The rectangle problem was still open in the 2015 second edition of Graham and Butler's \emph{Rudiments of Ramsey Theory} \cite{graham-butler}.

In 2026, Kova\v{c} answered the rectangle question by constructing a $25$-coloring of $\mathbb R^2$ with no monochromatic rectangle of area $1$ \cite{kovac}. For rectangles, scaling shows that it suffices to consider only rectangles of area 1, so this answers the question negatively. His coloring has a stronger property: it contains no monochromatic parallelogram whose two adjacent side lengths have product $1$. We improve the number of colors from $25$ to $21$.

\begin{theorem}\label{thm:main}
There exists a $21$-coloring of $\mathbb R^2$ with no monochromatic parallelogram whose two adjacent side lengths have product $1$.
\end{theorem}

For a rectangle, the product of two adjacent side lengths is its area, so the theorem gives a $21$-coloring with no monochromatic rectangle of area $1$. The analogous question for parallelograms of prescribed area is more tricky and remains open \cite{kovac}. If $u$ and $v$ are adjacent side vectors and $\theta$ is the angle between them, then the area is $|u||v||\sin\theta|$. The theorem rules out only the case $|u||v|=1$, and a parallelogram can have area $1$ with $|u||v|>1$; for instance, take $\sin\theta=1/(|u||v|)$.

Our construction follows Kova\v{c}'s approach. We identify the plane with $\mathbb C$, assign colors according to the location of $z^2$, and replace his periodic square tiling with the Voronoi tiling of the triangular lattice. The relevant sublattice has index $21$. Section~2 gives the coloring and proves Theorem~\ref{thm:main}; Section~3 shows that no smaller index works in this construction.

\section{Proof of Theorem~\ref{thm:main}}

Identify $\mathbb R^2$ with $\mathbb C$. Let $P=ABCD$ be a possibly degenerate parallelogram. We write its vertices as $z_A=z$, $z_B=z+u$, $z_C=z+u+v$, and $z_D=z+v$.


Consider the alternating square sum
\begin{equation}\label{eq:invariant}
\mathscr I(P)=z_A^2-z_B^2+z_C^2-z_D^2.
\end{equation}
Substituting the vertex coordinates gives $\mathscr I(P)=z^2-(z+u)^2+(z+u+v)^2-(z+v)^2=2uv$. Thus $|\mathscr I(P)|=2$ whenever the product of the two adjacent side lengths is $1$. We will therefore construct a coloring for which $\mathscr I(P)$ avoids the circle $|w|=2$ whenever $P$ is monochromatic.

Let $\omega=e^{\pi i/3}$ and $L=\mathbb Z+\mathbb Z\omega$. The Voronoi cell $H_0$ of $0$ in $L$ is a regular hexagon with inradius $1/2$ and circumradius $1/\sqrt3$. Choose a half-open version $V$ of $H_0$ by assigning each boundary point of the Voronoi tiling to one adjacent cell. Then the translates $\lambda+V$, with $\lambda\in L$, partition $\mathbb C$. Set $\alpha=4+\omega$.

For $j\in\{0,1,\ldots,20\}$, define
\[
\mathscr C_j=\bigl\{z\in\mathbb C:z^2\in \frac{5}{6}(p+q\omega+V)\text{ for some }p,q\in\mathbb Z\text{ with }p-4q\equiv j\pmod{21}\bigr\}.
\]
Because the cells $\frac{5}{6}(\lambda+V)$ partition $\mathbb C$, the sets $\mathscr C_0,\ldots,\mathscr C_{20}$ form a $21$-coloring of $\mathbb C$.

The congruence condition is equivalent to membership in $\alpha L$:
$p+q\omega\in\alpha L$ if and only if $p-4q\equiv0\pmod{21}$. If $p+q\omega=\alpha(r+s\omega)$, then $p=4r-s$ and $q=r+5s$, so $p-4q=-21s$. Conversely, if $p-4q=21k$, then $p+q\omega=\alpha((q+5k)-k\omega)$.

\begin{lemma}\label{lem:separation}
The set $\frac{5}{6}(\alpha L+4H_0)$ is disjoint from the circle $|w|=2$.
\end{lemma}

\begin{proof}
The hexagon $\frac{5}{6}(4H_0)$ has circumradius $10/(3\sqrt3)<2$, so the copy centered at $0$ lies inside the open disk $|w|<2$.

Let $0\ne h=p+q\omega\in L$. Since $|h|^2=p^2+pq+q^2$ and $p^2+pq+q^2\equiv(p-q)^2\pmod3$, the positive integer $|h|^2$ cannot equal $2$. Hence either $|h|=1$ or $|h|\ge\sqrt3$.

If $|h|\ge\sqrt3$, then $|\alpha h|\ge\sqrt{63}$, while $4H_0$ has circumradius $4/\sqrt3$. Hence $d(\alpha h,4H_0)\ge\sqrt{63}-4/\sqrt3>5/2$.
If $|h|=1$, multiplication by $h$ is a rotation preserving $L$ and $H_0$, so it is enough to consider $h=1$. We have $\alpha=9/2+(\sqrt3/2)i$, and the side of $4H_0$ on the line $\operatorname{Re}z=2$ has endpoints $2\pm2i/\sqrt3$. The orthogonal projection of $\alpha$ onto this side lies on the segment, so $d(\alpha,4H_0)=5/2$.

We have shown that $d(\alpha h,4H_0)\ge5/2$ for every nonzero $h\in L$. After scaling by $5/6$, this gives $d(0,\frac{5}{6}(\alpha h+4H_0))\ge25/12>2$. Thus every nonzero translate lies outside the disk $|w|\le2$, while the copy centered at $0$ lies inside it. None of the sets $\frac{5}{6}(\alpha h+4H_0)$ can therefore meet the circle $|w|=2$.
\end{proof}

\begin{proof}[Proof of Theorem~\ref{thm:main}]
Suppose that $P=ABCD$ is monochromatic, with all four vertices in $\mathscr C_j$. For each $X\in\{A,B,C,D\}$, write $z_X^2=\frac{5}{6}(\lambda_X+\varepsilon_X)$, where $\lambda_X=p_X+q_X\omega\in L$ satisfies $p_X-4q_X\equiv j\pmod{21}$ and $\varepsilon_X\in V\subset H_0$. The alternating sum $\lambda_A-\lambda_B+\lambda_C-\lambda_D$ is congruent to $j-j+j-j=0$ modulo $21$, so it lies in $\alpha L$. Since $H_0$ is centrally symmetric and convex, $\varepsilon_A-\varepsilon_B+\varepsilon_C-\varepsilon_D\in4H_0$.
It follows that $\mathscr I(P)\in\frac{5}{6}(\alpha L+4H_0)$. Lemma~\ref{lem:separation} shows that this set is disjoint from the circle $|w|=2$. Hence $|\mathscr I(P)|\ne2$, and the product of the two adjacent side lengths of $P$ is not $1$.
\end{proof}

\section{Sharpness of the lattice construction}

We next show that no smaller index can work in this construction. Let $M\subset L$ be a finite-index sublattice. Suppose that, after scaling, the copy of $4H_0$ centered at $0$ lies inside $|w|<2$, while every copy centered at a nonzero point of $M$ lies outside the closed disk $|w|\le2$. The first condition forces the scaling factor to be less than $\sqrt3/2$. The second can then hold only if $d(m,4H_0)>4/\sqrt3$ for every nonzero $m\in M$.

\begin{proposition}\label{prop:sharp}
If $M\subset L$ has index at most $20$, then there exists $0\ne m\in M$ such that $d(m,4H_0)\le4/\sqrt3$.
\end{proposition}

\begin{proof}
In lattice coordinates, a point $x+y\omega$ lies in $4H_0$ if and only if $|2x+y|\le4$, $|x+2y|\le4$, and $|y-x|\le4$. Let $Q=4H_0-(1+\omega)/3$. A lattice point $p+q\omega$ lies in $Q$ if and only if $|2p+q+1|\le4$, $|p+2q+1|\le4$, and $|q-p|\le4$.
The lattice points in $Q$ are:

\begin{center}
\begin{tabular}{c|l}
$q$ & $p$ \\ \hline
$-3$ & $1$ \\
$-2$ & $-1,0,1,2$ \\
$-1$ & $-2,-1,0,1,2$ \\
$0$ & $-2,-1,0,1$ \\
$1$ & $-3,-2,-1,0,1$ \\
$2$ & $-2,-1$
\end{tabular}
\end{center}

There are $21$ lattice points in $Q$. If $[L:M]\le20$, two of them, say $x$ and $y$, lie in the same coset of $M$. Their difference $m=x-y$ is a nonzero point of $M$. Since $Q$ is a translate of $4H_0$, we have $m\in Q-Q=8H_0$. Hence $m/2\in4H_0$ and $|m|\le8/\sqrt3$. It follows that $d(m,4H_0)\le|m-m/2|=|m|/2\le4/\sqrt3$.
\end{proof}

Multiplication by $\alpha=4+\omega$ scales Euclidean area by $|\alpha|^2=21$, so $[L:\alpha L]=21$. Lemma~\ref{lem:separation} shows that index $21$ works, while Proposition~\ref{prop:sharp} rules out every index at most $20$. Thus $21$ is the smallest index for which this sublattice construction can keep the translates away from the circle $|w|=2$.

\begin{remark}
  Recall that a coloring is Jordan-measurable when the boundary of each color class has Lebesgue measure zero. Kova\v{c}'s $25$-coloring is Jordan-measurable, and so is the coloring constructed here. For our coloring, the boundary of each color class is contained in the inverse image under $z\mapsto z^2$ of the union of the Voronoi-cell boundaries. This union is a countable union of line segments, and the inverse image of each segment is contained in a conic of planar measure zero.
\end{remark}

\begin{problem}\label{prob:min-colors}
Determine the least integer $r$ for which there exists an $r$-coloring of $\mathbb R^2$ with no monochromatic rectangle of area $1$.
\end{problem}

Theorem~\ref{thm:main} gives $r\le21$. Proposition~\ref{prop:sharp} shows that the construction of Section~2 cannot produce fewer colors.

\section*{Acknowledgments}

This work was completed as part of the REU on Combinatorics, AI, and Algorithms for Real Problems (REU-CAAR) at the University of Maryland. We thank Bill Gasarch for introducing us to this problem and for his guidance throughout this project.

{\small

\begin{thebibliography}{9}

\bibitem{erdos-graham}
P.~Erd\H{o}s and R.~L. Graham,
\emph{Old and new problems and results in combinatorial number theory},
Monographies de L'Enseignement Math\'ematique, vol.~28,
Universit\'e de Gen\`eve, L'Enseignement Math\'ematique,
Geneva, 1980.

\bibitem{graham}
R.~L. Graham,
\emph{On partitions of $E^n$},
J. Combin. Theory Ser. A 28 (1980), no.~1, 89--97,
\href{https://doi.org/10.1016/0097-3165(80)90061-8}
{doi:10.1016/0097-3165(80)90061-8}.

\bibitem{graham-butler}
R.~L. Graham and S.~Butler,
\emph{Rudiments of Ramsey Theory}, 2nd ed.,
CBMS Regional Conference Series in Mathematics, vol.~123,
Conference Board of the Mathematical Sciences, Washington, DC;
American Mathematical Society, Providence, RI, 2015,
\href{https://doi.org/10.1090/cbms/123}
{doi:10.1090/cbms/123}.

\bibitem{kovac}
V.~Kova\v{c},
\emph{Coloring and density theorems for configurations of a given volume},
Proc. Lond. Math. Soc. (3) 132 (2026), no.~3,
Paper No.~e70143, 56~pp.,
\href{https://doi.org/10.1112/plms.70143}
{doi:10.1112/plms.70143}.

\end{thebibliography}

}



\end{document}
